\section{Randomized --- Zero PRun}
%%% generated by the blueprint pipeline from TCSlib.Complexity.Randomized.ZeroPRun
\subsection*{Randomized --- Zero PRun: Introduction}

Declarations of the Lean module \texttt{TCSlib.Complexity.Randomized.ZeroPRun}.

\subsection*{Randomized --- Zero PRun: Results}

\begin{theorem}[Polynomial-Time Comparison of Output Lengths]
\lean{Randomized.PIT.polyTimeComputable_ltLen}
\label{Randomized.PIT.polyTimeComputable_ltLen}
\leanok
\uses{Complexity.PolyTimeComputable, Complexity.PolyTimeComputable.comp, Complexity.PolyTimeComputable.pairEncode, Complexity.pairFstD_pairEncode, Complexity.pairSndD_pairEncode, Complexity.polyTimeComputable_lenLe, Complexity.polyTimeComputable_not}
\difficulty{4}
Let $f,g : \{0,1\}^* \to \{0,1\}^*$ be polynomial-time computable functions. The
function that maps each bit string $D$ to the one-bit string whose bit is $1$ exactly
when $\abs{f(D)} < \abs{g(D)}$ is polynomial-time computable.
\end{theorem}

\begin{theorem}[Polynomial-Time Computability of a Tag Bit]
\lean{Randomized.PIT.polyTimeComputable_tagBit}
\label{Randomized.PIT.polyTimeComputable_tagBit}
\leanok
\uses{Complexity.BinArith.polyTimeComputable_drop, Complexity.PolyTimeComputable, Complexity.PolyTimeComputable.comp, Complexity.polyTimeComputable_headD, Randomized.PIT.tagBit}
\difficulty{1}
For every natural number $i$, the function sending a bit string $w$ to the one-bit word
consisting of the bit at position $i$ of $w$ (with positions counted from zero and
missing bits taken to be false) is polynomial-time computable.
\end{theorem}

\begin{theorem}[Polynomial-time computability of the word body]
\lean{Randomized.PIT.polyTimeComputable_bodyW}
\label{Randomized.PIT.polyTimeComputable_bodyW}
\leanok
\uses{Complexity.BinArith.polyTimeComputable_drop, Complexity.PolyTimeComputable, Randomized.PIT.bodyW}
\difficulty{1}
The function on bit strings that deletes the first four bits is polynomial-time
computable.
\end{theorem}

\begin{theorem}[Polynomial-time computability of the first field]
\lean{Randomized.PIT.polyTimeComputable_arg1W}
\label{Randomized.PIT.polyTimeComputable_arg1W}
\leanok
\uses{Complexity.BinArith.polyTimeComputable_uHead, Complexity.PolyTimeComputable, Complexity.PolyTimeComputable.comp, Randomized.PIT.arg1W, Randomized.PIT.polyTimeComputable_bodyW}
\difficulty{1}
The function on bit strings that extracts the first unary-delimited argument field from
the body of a word is polynomial-time computable.
\end{theorem}

\begin{theorem}[Polynomial-Time Computability of the Middle Word]
\lean{Randomized.PIT.polyTimeComputable_midW}
\label{Randomized.PIT.polyTimeComputable_midW}
\leanok
\uses{Complexity.BinArith.polyTimeComputable_uRest, Complexity.PolyTimeComputable, Complexity.PolyTimeComputable.comp, Randomized.PIT.midW, Randomized.PIT.polyTimeComputable_bodyW}
\difficulty{1}
The function that first deletes the first four bits of a Boolean word and then applies
the specified one-pass transducer is computable by a finite Turing machine in polynomial
time: there exist natural numbers $C$ and $c$ such that, on every input word of length
$n$, the machine halts within $C(n+1)^c$ steps and outputs the resulting word.
\end{theorem}

\begin{theorem}[Polynomial-time computability of the second field]
\lean{Randomized.PIT.polyTimeComputable_arg2W}
\label{Randomized.PIT.polyTimeComputable_arg2W}
\leanok
\uses{Complexity.BinArith.polyTimeComputable_uHead, Complexity.PolyTimeComputable, Complexity.PolyTimeComputable.comp, Randomized.PIT.arg2W, Randomized.PIT.polyTimeComputable_midW}
\difficulty{1}
The function on bit strings that extracts the second unary-delimited argument field from
the body of a word is polynomial-time computable.
\end{theorem}

\begin{theorem}[Polynomial-time computability of the second rest function]
\lean{Randomized.PIT.polyTimeComputable_rest2W}
\label{Randomized.PIT.polyTimeComputable_rest2W}
\leanok
\uses{Complexity.BinArith.polyTimeComputable_uRest, Complexity.PolyTimeComputable, Complexity.PolyTimeComputable.comp, Randomized.PIT.polyTimeComputable_midW, Randomized.PIT.rest2W}
\difficulty{1}
The function that maps a bit string $w$ to the result of applying the unary rest
transducer to the result of applying that transducer to the suffix of $w$ obtained by
deleting its first four bits is computable by a finite Turing machine in polynomial
time.
\end{theorem}

\begin{theorem}[Polynomial-time computability of a bit predicate]
\lean{Randomized.PIT.polyTimeComputable_isVarW}
\label{Randomized.PIT.polyTimeComputable_isVarW}
\leanok
\uses{Complexity.PolyTimeComputable, Complexity.polyTimeComputable_and, Complexity.polyTimeComputable_not, Randomized.PIT.isVarW, Randomized.PIT.polyTimeComputable_tagBit}
\difficulty{3}
For every finite bit string $w$, the function that outputs the one-bit string $[b]$,
where $b$ is true exactly when the bits of $w$ at positions $1$, $2$, and $3$ are all
false (with a missing bit treated as false), is polynomial-time computable.
\end{theorem}

\begin{theorem}[Polynomial-Time Unit-Tag Recognition]
\lean{Randomized.PIT.polyTimeComputable_isOneW}
\label{Randomized.PIT.polyTimeComputable_isOneW}
\leanok
\uses{Complexity.PolyTimeComputable, Complexity.polyTimeComputable_and, Complexity.polyTimeComputable_not, Randomized.PIT.isOneW, Randomized.PIT.polyTimeComputable_tagBit}
\difficulty{3}
The function on bit strings that returns the one-bit word $[1]$ when the second and
third bits are $0$ and the fourth bit is $1$, and returns $[0]$ otherwise, is
polynomial-time computable. A deterministic Turing machine with finitely many states
computes it on every input of length $n$ within $C(n+1)^c$ steps for some
$C,c\in\mathbb{N}$.
\end{theorem}

\begin{theorem}[Polynomial-Time Addition-Tag Recognition]
\lean{Randomized.PIT.polyTimeComputable_isAddW}
\label{Randomized.PIT.polyTimeComputable_isAddW}
\leanok
\uses{Complexity.PolyTimeComputable, Complexity.polyTimeComputable_and, Complexity.polyTimeComputable_not, Randomized.PIT.isAddW, Randomized.PIT.polyTimeComputable_tagBit}
\difficulty{3}
The function on bit strings that returns the one-bit word $[1]$ when the second and
fourth bits are $1$ and the third bit is $0$, and returns $[0]$ otherwise, is
polynomial-time computable. In particular, a deterministic Turing machine with finitely
many states computes it on every input of length $n$ within $C(n+1)^c$ steps for some
$C,c\in\mathbb{N}$.
\end{theorem}

\begin{theorem}[Polynomial-Time Subtraction-Tag Recognition]
\lean{Randomized.PIT.polyTimeComputable_isSubW}
\label{Randomized.PIT.polyTimeComputable_isSubW}
\leanok
\uses{Complexity.PolyTimeComputable, Complexity.polyTimeComputable_and, Complexity.polyTimeComputable_not, Randomized.PIT.isSubW, Randomized.PIT.polyTimeComputable_tagBit}
\difficulty{3}
The function on bit strings that returns the one-bit word $[1]$ when the second, third,
and fourth bits are all $1$, and returns $[0]$ otherwise, is polynomial-time computable.
A deterministic Turing machine with finitely many states computes it on every input of
length $n$ within $C(n+1)^c$ steps for some $C,c\in\mathbb{N}$.
\end{theorem}

\begin{theorem}[Polynomial-Time Multiplication-Tag Recognition]
\lean{Randomized.PIT.polyTimeComputable_isMulW}
\label{Randomized.PIT.polyTimeComputable_isMulW}
\leanok
\uses{Complexity.PolyTimeComputable, Complexity.polyTimeComputable_and, Complexity.polyTimeComputable_not, Randomized.PIT.isMulW, Randomized.PIT.polyTimeComputable_tagBit}
\difficulty{3}
The function on bit strings that returns the one-bit word $[1]$ when the second bit is
$1$ and the third and fourth bits are $0$, and returns $[0]$ otherwise, is
polynomial-time computable. A deterministic Turing machine with finitely many states
computes it on every input of length $n$ within $C(n+1)^c$ steps for some
$C,c\in\mathbb{N}$.
\end{theorem}

\begin{theorem}[Polynomial-Time Completion Recognition]
\lean{Randomized.PIT.polyTimeComputable_isDone}
\label{Randomized.PIT.polyTimeComputable_isDone}
\leanok
\uses{Complexity.BinArith.polyTimeComputable_drop, Complexity.PolyTimeComputable, Complexity.PolyTimeComputable.comp, Complexity.polyTimeComputable_and, Complexity.polyTimeComputable_headD, Complexity.polyTimeComputable_isNil, Complexity.polyTimeComputable_not, Complexity.polyTimeComputable_or, Randomized.PIT.isDone}
\difficulty{3}
The function on bit strings that returns $[1]$ exactly when the input is empty or
consists of the single bit $0$, and returns $[0]$ otherwise, is polynomial-time
computable. A deterministic Turing machine with finitely many states computes it on
every input of length $n$ within $C(n+1)^c$ steps for some $C,c\in\mathbb{N}$.
\end{theorem}

\begin{theorem}[Polynomial-Time Computability of the Rest Transformation]
\lean{Randomized.PIT.polyTimeComputable_restW}
\label{Randomized.PIT.polyTimeComputable_restW}
\leanok
\uses{Complexity.PolyTimeComputable, Complexity.polyTimeComputable_ite, Randomized.PIT.polyTimeComputable_bodyW, Randomized.PIT.polyTimeComputable_isOneW, Randomized.PIT.polyTimeComputable_isVarW, Randomized.PIT.polyTimeComputable_midW, Randomized.PIT.polyTimeComputable_rest2W, Randomized.PIT.restW}
\difficulty{3}
The function on finite bit strings that sends each string $w$ to the word obtained by
the defining case distinction for the rest transformation is computable by a
deterministic Turing machine in time bounded by a polynomial in the input length.
\end{theorem}

\begin{theorem}[Polynomial-Time Computability of a Nested Projection]
\lean{Randomized.PIT.polyTimeComputable_sW}
\label{Randomized.PIT.polyTimeComputable_sW}
\leanok
\uses{Complexity.BinArith.polyTimeComputable_slot2, Complexity.PolyTimeComputable, Randomized.PIT.sW}
\difficulty{1}
The function on bit strings obtained by applying the total second projection twice and
then the total first projection is polynomial-time computable.
\end{theorem}

\begin{theorem}[Polynomial-Time Computability of the Input Projection]
\lean{Randomized.PIT.polyTimeComputable_sI}
\label{Randomized.PIT.polyTimeComputable_sI}
\leanok
\uses{Complexity.BinArith.polyTimeComputable_slot3, Complexity.PolyTimeComputable, Complexity.PolyTimeComputable.comp, Complexity.polyTimeComputable_pairFstD, Randomized.PIT.sI}
\difficulty{2}
The transformation that sends each finite bit string $D$ to $\mathrm{sI}(D)$ is
polynomial-time computable.
\end{theorem}

\begin{theorem}[Polynomial-Time Computability of the Length Projection]
\lean{Randomized.PIT.polyTimeComputable_sL}
\label{Randomized.PIT.polyTimeComputable_sL}
\leanok
\uses{Complexity.BinArith.polyTimeComputable_slot4, Complexity.PolyTimeComputable, Randomized.PIT.sL}
\difficulty{1}
The transformation that sends each finite bit string $D$ to $\mathrm{sL}(D)$ is
polynomial-time computable.
\end{theorem}

\begin{theorem}[Polynomial-Time Computability of the Context Projection]
\lean{Randomized.PIT.polyTimeComputable_sCtx}
\label{Randomized.PIT.polyTimeComputable_sCtx}
\leanok
\uses{Complexity.PolyTimeComputable, Complexity.polyTimeComputable_pairFstD, Randomized.PIT.sCtx}
\difficulty{1}
The transformation that sends each finite bit string $D$ to $\mathrm{sCtx}(D)$ is
polynomial-time computable.
\end{theorem}

\begin{theorem}[Polynomial-Time Computability of the Number Projection]
\lean{Randomized.PIT.polyTimeComputable_sN}
\label{Randomized.PIT.polyTimeComputable_sN}
\leanok
\uses{Complexity.PolyTimeComputable, Complexity.PolyTimeComputable.comp, Complexity.polyTimeComputable_pairFstD, Randomized.PIT.polyTimeComputable_sCtx, Randomized.PIT.sN}
\difficulty{1}
The transformation that sends each finite bit string $D$ to $\mathrm{sN}(D)$ is
polynomial-time computable.
\end{theorem}

\begin{theorem}[Polynomial-Time Computability of the Key Transformation]
\lean{Randomized.PIT.polyTimeComputable_sK}
\label{Randomized.PIT.polyTimeComputable_sK}
\leanok
\uses{Complexity.BinArith.polyTimeComputable_add, Complexity.BinArith.polyTimeComputable_slot2, Complexity.PolyTimeComputable, Complexity.PolyTimeComputable.comp, Complexity.PolyTimeComputable.pairEncode, Complexity.polyTimeComputable_const, Randomized.PIT.polyTimeComputable_sCtx, Randomized.PIT.sK}
\difficulty{1}
The transformation that sends each finite bit string $D$ to $\mathrm{sK}(D)$ is
polynomial-time computable.
\end{theorem}

\begin{theorem}[Polynomial-Time Computability of Indexed Circuit Values]
\lean{Randomized.PIT.polyTimeComputable_xW}
\label{Randomized.PIT.polyTimeComputable_xW}
\leanok
\uses{Complexity.BinArith.polyTimeComputable_getW, Complexity.PolyTimeComputable, Complexity.PolyTimeComputable.comp, Complexity.PolyTimeComputable.pairEncode, Randomized.PIT.polyTimeComputable_sL, Randomized.PIT.xW}
\difficulty{1}
For every polynomial-time computable function $f$ on bit strings, the function sending
an encoded gate-state string $D$ to the value obtained by using $f(D)$ as the index and
$D$ as the state is polynomial-time computable.
\end{theorem}

\begin{theorem}[Polynomial-Time Computability of Modular Gate Evaluation]
\lean{Randomized.PIT.polyTimeComputable_valW}
\label{Randomized.PIT.polyTimeComputable_valW}
\leanok
\uses{Complexity.BinArith.getW, Complexity.BinArith.polyTimeComputable_add, Complexity.BinArith.polyTimeComputable_getW, Complexity.BinArith.polyTimeComputable_mod, Complexity.BinArith.polyTimeComputable_mul, Complexity.BinArith.polyTimeComputable_slot3, Complexity.BinArith.polyTimeComputable_sub, Complexity.PolyTimeComputable, Complexity.PolyTimeComputable.comp, Complexity.PolyTimeComputable.pairEncode, Complexity.polyTimeComputable_const, Complexity.polyTimeComputable_ite, Randomized.PIT.arg1W, Randomized.PIT.arg2W, Randomized.PIT.polyTimeComputable_arg1W, Randomized.PIT.polyTimeComputable_arg2W, Randomized.PIT.polyTimeComputable_isAddW, Randomized.PIT.polyTimeComputable_isOneW, Randomized.PIT.polyTimeComputable_isSubW, Randomized.PIT.polyTimeComputable_isVarW, Randomized.PIT.polyTimeComputable_sCtx, Randomized.PIT.polyTimeComputable_sK, Randomized.PIT.polyTimeComputable_sW, Randomized.PIT.polyTimeComputable_xW, Randomized.PIT.rawValW, Randomized.PIT.sPts, Randomized.PIT.sW, Randomized.PIT.valW, Randomized.PIT.xW, Turing.pairEncode}
\difficulty{4}
The function sending each encoded gate-state string $D$ to its modular value is
computable by a deterministic Turing machine in time bounded by a polynomial in the
input length.
\end{theorem}

\begin{theorem}[Polynomial-time computability of the word-validity test]
\lean{Randomized.PIT.polyTimeComputable_okW}
\label{Randomized.PIT.polyTimeComputable_okW}
\leanok
\uses{Complexity.BinArith.hasFalse, Complexity.BinArith.polyTimeComputable_hasFalse, Complexity.PolyTimeComputable, Complexity.PolyTimeComputable.comp, Complexity.polyTimeComputable_and, Complexity.polyTimeComputable_headD, Complexity.polyTimeComputable_or, Randomized.PIT.arg1W, Randomized.PIT.arg2W, Randomized.PIT.bodyW, Randomized.PIT.isAddW, Randomized.PIT.isMulW, Randomized.PIT.isSubW, Randomized.PIT.isVarW, Randomized.PIT.midW, Randomized.PIT.okW, Randomized.PIT.polyTimeComputable_arg1W, Randomized.PIT.polyTimeComputable_arg2W, Randomized.PIT.polyTimeComputable_bodyW, Randomized.PIT.polyTimeComputable_isAddW, Randomized.PIT.polyTimeComputable_isMulW, Randomized.PIT.polyTimeComputable_isOneW, Randomized.PIT.polyTimeComputable_isSubW, Randomized.PIT.polyTimeComputable_isVarW, Randomized.PIT.polyTimeComputable_ltLen, Randomized.PIT.polyTimeComputable_midW, Randomized.PIT.polyTimeComputable_sI, Randomized.PIT.polyTimeComputable_sN, Randomized.PIT.polyTimeComputable_sW, Randomized.PIT.sI, Randomized.PIT.sN, Randomized.PIT.sW}
\difficulty{4}
The Boolean-valued function that maps a bit string $D$ to the one-bit string whose bit
is the result of the following test is computable in polynomial time: extract
$s_N,s_I,s_W$ by the successive total first-projection operations specified above, let
$n$ and $i$ be the lengths of $s_N$ and $s_I$, respectively, and test whether $s_W$
begins with $1$ and satisfies at least one of these conditions: it is tagged as a
variable, its body contains a $0$, and the first argument extracted from the body has
length less than $n$; it is tagged as the constant $1$; or it is tagged as addition,
subtraction, or multiplication, its body and middle word each contain a $0$, and both
extracted argument lengths are less than $i$.
\end{theorem}

\begin{theorem}[Polynomial-Time Gate-State Update]
\lean{Randomized.PIT.polyTimeComputable_gStep}
\label{Randomized.PIT.polyTimeComputable_gStep}
\leanok
\uses{Complexity.PolyTimeComputable, Complexity.PolyTimeComputable.append, Complexity.PolyTimeComputable.comp, Complexity.PolyTimeComputable.pairEncode, Complexity.polyTimeComputable_const, Complexity.polyTimeComputable_id, Complexity.polyTimeComputable_ite, Complexity.polyTimeComputable_prepend, Randomized.PIT.gStep, Randomized.PIT.polyTimeComputable_isDone, Randomized.PIT.polyTimeComputable_okW, Randomized.PIT.polyTimeComputable_restW, Randomized.PIT.polyTimeComputable_sCtx, Randomized.PIT.polyTimeComputable_sI, Randomized.PIT.polyTimeComputable_sL, Randomized.PIT.polyTimeComputable_sW, Randomized.PIT.polyTimeComputable_valW}
\difficulty{3}
The bit-string update that leaves a state unchanged when its work word is done,
otherwise checks whether that word is valid, and then either advances the computation by
updating the encoded work word, counter, and value list or resets the work word as
specified, is polynomial-time computable. That is, a deterministic Turing machine with
finitely many states computes this update on every bit-string input in at most
$C(n+1)^c$ steps for some $C,c\in\mathbb{N}$, where $n$ is the input length.
\end{theorem}

\begin{theorem}[First projection preserved by a step]
\lean{Randomized.PIT.pairFstD_gStep}
\label{Randomized.PIT.pairFstD_gStep}
\leanok
\uses{Complexity.pairFstD, Randomized.PIT.gStep, Randomized.PIT.sCtx}
\difficulty{3}
For every bit string $D$, the first component extracted from the result of one
gate-processing step on $D$ equals the first component extracted from $D$.
\end{theorem}

\begin{theorem}[Length bound for the remaining word]
\lean{Randomized.PIT.length_restW_le}
\label{Randomized.PIT.length_restW_le}
\leanok
\uses{Complexity.BinArith.length_uRest_le, Randomized.PIT.bodyW, Randomized.PIT.midW, Randomized.PIT.rest2W, Randomized.PIT.restW}
\difficulty{4}
For every bit string $w$, the word selected as the remainder according to whether $w$ is
tagged as a variable, the constant one, or another operation has length at most $|w|$.
\end{theorem}

\begin{theorem}[Length bound for the reduced value]
\lean{Randomized.PIT.length_valW_le}
\label{Randomized.PIT.length_valW_le}
\leanok
\uses{Complexity.BinArith.add, Complexity.BinArith.length_bits_le_of_le, Complexity.BinArith.mod, Complexity.LogProg.bitsVal, Complexity.LogProg.bitsVal_bits, Complexity.pairFstD_pairEncode, Complexity.pairSndD_pairEncode, Randomized.PIT.rawValW, Randomized.PIT.sK, Randomized.PIT.sRaw, Randomized.PIT.valW}
\difficulty{4}
For every bit string $D$, the word representing the value reduced modulo the key
extracted from $D$ has length at most the length of the raw input extracted from $D$.
\end{theorem}

\begin{theorem}[Length bound for one update]
\lean{Randomized.PIT.length_gStep_le}
\label{Randomized.PIT.length_gStep_le}
\leanok
\uses{Complexity.length_pairFstD_le, Complexity.length_pairSndD_le, Complexity.length_pair_components_le, Complexity.pairFstD, Complexity.pairSndD, Randomized.PIT.gStep, Randomized.PIT.length_restW_le, Randomized.PIT.length_valW_le, Randomized.PIT.sCtx, Randomized.PIT.sI, Randomized.PIT.sL, Randomized.PIT.sRaw, Randomized.PIT.sW, Turing.length_pairEncode}
\difficulty{5}
For every bit string $D$, let $T(D)$ be the next state produced by the update rule that
tests whether the current word is done or valid and, if it is valid, advances the index
and appends the computed value. Then $|T(D)|\le |D|+10(|\pi_1(D)|+1)$, where $\pi_1$ is
the total first projection of the pair encoding.
\end{theorem}

\begin{definition}[Gate initialization]
\lean{Randomized.PIT.gateInit}
\label{Randomized.PIT.gateInit}
\leanok
\uses{Complexity.pairFstD, Complexity.pairSndD, Turing.pairEncode}
For a Boolean list $z$, initialize the gate state by pairing the second component
extracted from $z$ with a nested encoding: the first component extracted from $z$,
paired with the second component extracted from $z$, paired with the encoding of two
empty lists.
\end{definition}

\begin{definition}[Gate execution]
\lean{Randomized.PIT.gateRun}
\label{Randomized.PIT.gateRun}
\leanok
\uses{Complexity.BinArith.forLoop, Complexity.pairSndD, Randomized.PIT.gStep, Randomized.PIT.gateInit}
For a Boolean list $z$, run the loop that repeatedly applies the gate step to the
initialized state, then return the second component extracted from the resulting state.
\end{definition}

\begin{theorem}[Polynomial-Time Gate Evaluation]
\lean{Randomized.PIT.polyTimeComputable_gateRun}
\label{Randomized.PIT.polyTimeComputable_gateRun}
\leanok
\uses{Complexity.BinArith.forLoop, Complexity.BinArith.forLoop_bound_of_ctx, Complexity.BinArith.polyTimeComputable_forLoop, Complexity.PolyTimeComputable, Complexity.PolyTimeComputable.comp, Complexity.PolyTimeComputable.pairEncode, Complexity.polyTimeComputable_const, Complexity.polyTimeComputable_pairFstD, Complexity.polyTimeComputable_pairSndD, Randomized.PIT.gStep, Randomized.PIT.gateInit, Randomized.PIT.gateRun, Randomized.PIT.length_gStep_le, Randomized.PIT.pairFstD_gStep, Randomized.PIT.polyTimeComputable_gStep}
\difficulty{3}
The function that initializes the encoded gate-evaluation state from an input bit
string, iterates the gate-state update for the prescribed number of rounds, and returns
the second component of the resulting state is polynomial-time computable. Thus a
deterministic Turing machine with finitely many states computes it on every input of
length $n$ within $C(n+1)^c$ steps for some $C,c\in\mathbb{N}$.
\end{theorem}

\begin{theorem}[Encoded gate evaluation]
\lean{Randomized.PIT.gateRun_eq}
\label{Randomized.PIT.gateRun_eq}
\leanok
\uses{Complexity.BinArith.forLoop_pairEncode, Complexity.LogProg.bitsVal, Complexity.pairSndD_pairEncode, Randomized.PIT.absStep, Randomized.PIT.ctxOf, Randomized.PIT.encState, Randomized.PIT.gateInit, Randomized.PIT.gateRun, Randomized.PIT.iterate_gStep, Randomized.PIT.ptVal, Turing.pairEncode}
\difficulty{3}
Let $n$ be a natural number, let $R$ and $G$ be finite bit strings, and let $P$ be a
finite list of finite bit strings. The gate-running procedure applied to the encoding of
the context determined by $n$, $R$, and $P$, paired with $G$, returns the encoding of
that context together with the word and list of values obtained by applying the abstract
evaluation step $|G|$ times, starting from $(G,[])$, using modulus $\operatorname{bitsVal}(R)+1$ and
assigning the $i$th input variable the value $\operatorname{bitsVal}(P_i)+1$ (with an empty word when
$i$ is outside the list).
\end{theorem}

\begin{theorem}[Evaluation of an encoded gate list]
\lean{Randomized.PIT.gateRun_encodeGates}
\label{Randomized.PIT.gateRun_encodeGates}
\leanok
\uses{Complexity.LogProg.bitsVal, Randomized.AlgCircuit.encodeGates, Randomized.AlgCircuit.length_encodeGates, Randomized.AlgGate, Randomized.AlgGate.refsBelow, Randomized.PIT.ctxOf, Randomized.PIT.encState, Randomized.PIT.evalVals, Randomized.PIT.gateRun, Randomized.PIT.gateRun_eq, Randomized.PIT.iterate_absStep_encodeGates, Randomized.PIT.ptVal, Turing.pairEncode}
\difficulty{4}
Let $n$ be a natural number, let $R$ be a finite bit string, let $P$ be a finite list of
finite bit strings, and let $gs$ be a finite list of arithmetic gates with input
variables among $n$ variables. Assume that each gate at position $i$ that is an
addition, subtraction, or multiplication refers only to gate positions strictly less
than $i$. Put $m=\mathrm{val}(R)+1$, and assign input variable $j$ the value
$\mathrm{val}(P_j)+1$, taking the value to be $0$ when $j$ is beyond the end of $P$.
Then running the gate evaluator on the encoding of $gs$, paired with the context
determined by $n$, $R$, and $P$, produces the state encoding with that same context,
terminal word $[0]$, and the list of gate values evaluated successively modulo $m$,
starting from the empty list of prior values.
\end{theorem}

\begin{definition}[Successful run predicate]
\lean{Randomized.PIT.runOk}
\label{Randomized.PIT.runOk}
\leanok
\uses{Complexity.isNilB, Randomized.PIT.isDone, Randomized.PIT.sW}
A Boolean list $D$ is deemed to have a successful run exactly when the result of
applying $sW$ to $D$ is nonempty and satisfies the completion predicate.
\end{definition}

\begin{definition}[Run value extraction]
\lean{Randomized.PIT.runVal}
\label{Randomized.PIT.runVal}
\leanok
\uses{Complexity.BinArith.getW, Randomized.PIT.sI, Randomized.PIT.sL, Turing.pairEncode}
For a Boolean list $D$, extract the value by applying $getW$ to the encoding of the list
obtained from $sI(D)$ by deleting its first entry together with $sL(D)$.
\end{definition}

\begin{theorem}[Run Validity for an Encoded State]
\lean{Randomized.PIT.runOk_encState}
\label{Randomized.PIT.runOk_encState}
\leanok
\uses{Complexity.isNilB, Complexity.pairFstD_pairEncode, Complexity.pairSndD_pairEncode, Randomized.PIT.encState, Randomized.PIT.isDone, Randomized.PIT.runOk, Randomized.PIT.sW}
\difficulty{4}
For any bit strings $ctx,w$ and any finite list $vs$ of natural numbers, the
run-validity predicate of the state encoding formed from $ctx$, $w$, and $vs$ equals
true exactly when $w$ is the one-bit string $[0]$.
\end{theorem}

\begin{theorem}[Value of an Encoded State Run]
\lean{Randomized.PIT.runVal_encState}
\label{Randomized.PIT.runVal_encState}
\leanok
\uses{Complexity.BinArith.getW_pairEncode_bits, Complexity.pairFstD_pairEncode, Complexity.pairSndD_pairEncode, Randomized.PIT.encState, Randomized.PIT.runVal, Randomized.PIT.sI, Randomized.PIT.sL}
\difficulty{3}
For any bit strings $ctx,w$ and any finite list $vs$ of natural numbers, running the
value extraction on their encoded state yields the binary expansion of the entry of $vs$
at index $|vs|-1$, with index subtraction truncated at zero and value $0$ used when that
index is out of range.
\end{theorem}

\begin{theorem}[Accepted Circuit Encoding and Evaluation]
\lean{Randomized.PIT.exists_of_runOk}
\label{Randomized.PIT.exists_of_runOk}
\leanok
\uses{Complexity.LogProg.bitsVal, Randomized.AlgCircuit.encodeGates, Randomized.AlgGate, Randomized.AlgGate.refsBelow, Randomized.PIT.absStep, Randomized.PIT.ctxOf, Randomized.PIT.encState, Randomized.PIT.evalVals, Randomized.PIT.exists_of_iterate_absStep, Randomized.PIT.gateRun, Randomized.PIT.gateRun_eq, Randomized.PIT.ptVal, Randomized.PIT.runOk, Randomized.PIT.runOk_encState, Turing.pairEncode}
\difficulty{5}
Let $n$ be a natural number, let $R$ be a finite bit string, let $P$ be a finite list of
finite bit strings, and let $G$ be a finite bit string. If the run initialized from the
context determined by $n$, $R$, and $P$, together with $G$, has a nonempty word that is
done, then there is a finite list $gs$ of arithmetic gates on $n$ inputs such that $G$
is the encoding of $gs$, every gate at position $i$ refers only to gates at positions
less than $i$, and the run equals the encoding of that context, the word consisting of
the single bit $0$, and the list of gate values obtained by evaluating $gs$ modulo
$\mathrm{bitsVal}(R)+1$ with input values $i \mapsto \mathrm{bitsVal}(P_i)+1$.
\end{theorem}

\begin{theorem}[Head transducer output is unary]
\lean{Randomized.PIT.uHead_eq_replicate}
\label{Randomized.PIT.uHead_eq_replicate}
\leanok
\uses{Complexity.BinArith.uHead}
\difficulty{4}
For every bit string $w$, the output of the head transducer on $w$ is a string
consisting entirely of $1$ bits, of length equal to that output's length.
\end{theorem}

\begin{definition}[Scan execution]
\lean{Randomized.PIT.scanRun}
\label{Randomized.PIT.scanRun}
\leanok
\uses{Complexity.BinArith.uHead, Complexity.BinArith.uRest, Randomized.PIT.gateRun, Turing.pairEncode}
For a Boolean list $x$, run the gate procedure on the encoding whose first component is
the encoding of $uHead(x)$ with two empty lists and whose second component is
$uRest(x)$.
\end{definition}

\begin{definition}[Validity predicate]
\lean{Randomized.PIT.validW}
\label{Randomized.PIT.validW}
\leanok
\uses{Complexity.BinArith.hasFalse, Complexity.isNilB, Randomized.PIT.runOk, Randomized.PIT.sI, Randomized.PIT.scanRun}
A Boolean list $x$ is valid exactly when it contains a false entry, the scan run on $x$
is successful, and applying $sI$ to that scan run gives a nonempty result.
\end{definition}

\begin{theorem}[Unary scan context]
\lean{Randomized.PIT.scan_ctx}
\label{Randomized.PIT.scan_ctx}
\leanok
\uses{Complexity.BinArith.uHead, Randomized.PIT.ctxOf, Randomized.PIT.uHead_eq_replicate, Turing.pairEncode}
\difficulty{3}
For every bit string $x$, let $H$ be the output of the head transducer started in its
initial state on input $x$. Pairing $H$ with the empty string gives the same encoding as
the context whose unary prefix has length $|H|$ and whose other two components are
empty.
\end{theorem}

\begin{theorem}[Characterization of Valid Circuit Encodings]
\lean{Randomized.PIT.validW_iff}
\label{Randomized.PIT.validW_iff}
\leanok
\uses{Complexity.BinArith.eq_encodeNat_of_hasFalse, Complexity.BinArith.hasFalse_encodeNat, Complexity.BinArith.uHead, Complexity.BinArith.uHead_encodeNat, Complexity.BinArith.uRest_encodeNat, Complexity.isNilB, Complexity.pairFstD_pairEncode, Complexity.pairSndD_pairEncode, Randomized.AlgCircuit, Randomized.AlgCircuit.WellFormed, Randomized.AlgCircuit.decode, Randomized.AlgCircuit.decode_eq_some_iff, Randomized.AlgCircuit.encode, Randomized.AlgCircuit.encode_of_decode, Randomized.PIT.ctxOf, Randomized.PIT.encState, Randomized.PIT.exists_of_runOk, Randomized.PIT.gateRun_encodeGates, Randomized.PIT.length_evalVals, Randomized.PIT.runOk_encState, Randomized.PIT.sI, Randomized.PIT.scanRun, Randomized.PIT.scan_ctx, Randomized.PIT.validW, Turing.pairEncode}
\difficulty{5}
For every finite bit string $x$, the validity test returns true exactly when there is an
algebraic circuit $C$ on some finite number of inputs such that decoding $x$ yields $C$
and $C$ is well formed.
\end{theorem}

\begin{theorem}[Scan Run Gate Count]
\lean{Randomized.PIT.sI_scanRun}
\label{Randomized.PIT.sI_scanRun}
\leanok
\uses{Complexity.BinArith.uHead_encodeNat, Complexity.BinArith.uRest_encodeNat, Randomized.AlgCircuit, Randomized.AlgCircuit.WellFormed, Randomized.AlgCircuit.decode, Randomized.AlgCircuit.encode, Randomized.AlgCircuit.encode_of_decode, Randomized.AlgCircuit.size, Randomized.PIT.ctxOf, Randomized.PIT.encState, Randomized.PIT.gateRun_encodeGates, Randomized.PIT.length_evalVals, Randomized.PIT.sI, Randomized.PIT.scanRun, Turing.pairEncode}
\difficulty{3}
Let $x$ be a bit string, let $n$ be a natural number, and let $C$ be an algebraic
circuit on $n$ variables. If decoding $x$ yields exactly the pair $(n,C)$ and $C$ is
well formed, then the index component of the scan run on $x$ is the all-ones bit string
of length equal to the number of gates of $C$.
\end{theorem}

\begin{definition}[Evaluation modulo the scan value]
\lean{Randomized.PIT.evalModW}
\label{Randomized.PIT.evalModW}
\leanok
\uses{Randomized.PIT.gateRun, Randomized.PIT.runVal}
For a Boolean list $z$, evaluate it by extracting the run value from the gate execution
on $z$.
\end{definition}

\begin{theorem}[Correctness of modular circuit evaluation]
\lean{Randomized.PIT.evalModW_spec}
\label{Randomized.PIT.evalModW_spec}
\leanok
\uses{Complexity.LogProg.bitsVal, Complexity.LogProg.bitsVal_bits, Randomized.AlgCircuit, Randomized.AlgCircuit.WellFormed, Randomized.AlgCircuit.encodeGates, Randomized.AlgCircuit.eval, Randomized.PIT.ctxOf, Randomized.PIT.evalModW, Randomized.PIT.evalVals, Randomized.PIT.gateRun_encodeGates, Randomized.PIT.getD_map_natCast, Randomized.PIT.length_evalVals, Randomized.PIT.lt_of_mem_evalVals, Randomized.PIT.natCast_evalVals, Randomized.PIT.ptVal, Randomized.PIT.runVal_encState, Turing.pairEncode}
\difficulty{4}
Let $n$ be a natural number, let $R$ be a finite bit string, let $P$ be a finite list of
finite bit strings, and let $C$ be an arithmetic circuit with $n$ inputs. Write $b(w)$
for the natural number represented by a bit string $w$, with its first bit least
significant, and suppose $C$ has at least one gate and every gate that refers to earlier
gates refers only to gates preceding it in the list. The bit string produced by the
encoded evaluator on the encoding of the context and gate list has value less than
$b(R)+1$, and its residue modulo $b(R)+1$ equals the value of $C$ over
$\mathbb{Z}/(b(R)+1)\mathbb{Z}$ at the input whose $i$th coordinate is the residue of
$b(P_i)+1$.
\end{theorem}

\begin{theorem}[Polynomial-Time Computability of the Scan Run]
\lean{Randomized.PIT.polyTimeComputable_scanRun}
\label{Randomized.PIT.polyTimeComputable_scanRun}
\leanok
\uses{Complexity.BinArith.polyTimeComputable_uHead, Complexity.BinArith.polyTimeComputable_uRest, Complexity.BinArith.uHead, Complexity.BinArith.uRest, Complexity.PolyTimeComputable, Complexity.PolyTimeComputable.comp, Complexity.PolyTimeComputable.pairEncode, Complexity.polyTimeComputable_const, Randomized.PIT.gateRun, Randomized.PIT.polyTimeComputable_gateRun, Randomized.PIT.scanRun, Turing.pairEncode}
\difficulty{3}
For a bit string $x$, form the self-delimiting pair whose first component is the pair of
the unary head of $x$ and the pair of two empty strings, and whose second component is
the unary rest of $x$. Apply the gate-run function to this encoding. The resulting
function of $x$ is polynomial-time computable.
\end{theorem}

\begin{theorem}[Polynomial-Time Computability of the Run-Validity Test]
\lean{Randomized.PIT.polyTimeComputable_runOk}
\label{Randomized.PIT.polyTimeComputable_runOk}
\leanok
\uses{Complexity.PolyTimeComputable, Complexity.PolyTimeComputable.comp, Complexity.isNilB, Complexity.polyTimeComputable_and, Complexity.polyTimeComputable_isNil, Complexity.polyTimeComputable_not, Randomized.PIT.isDone, Randomized.PIT.polyTimeComputable_isDone, Randomized.PIT.polyTimeComputable_sW, Randomized.PIT.runOk, Randomized.PIT.sW}
\difficulty{3}
There exist a finite Turing machine $M$ over the binary alphabet and natural numbers
$C,c$ such that, on every finite bit string $D$, the machine halts within $C(|D|+1)^c$
steps and outputs the one-bit string whose bit is true exactly when the second component
obtained by twice projecting $D$ onto the second component and then onto the first
component is nonempty and is either empty or begins with $0$ followed by an empty
suffix.
\end{theorem}

\begin{theorem}[Polynomial-Time Computability of Encoding Validity]
\lean{Randomized.PIT.polyTimeComputable_validW}
\label{Randomized.PIT.polyTimeComputable_validW}
\leanok
\uses{Complexity.BinArith.hasFalse, Complexity.BinArith.polyTimeComputable_hasFalse, Complexity.PolyTimeComputable, Complexity.PolyTimeComputable.comp, Complexity.isNilB, Complexity.polyTimeComputable_and, Complexity.polyTimeComputable_isNil, Complexity.polyTimeComputable_not, Randomized.PIT.polyTimeComputable_runOk, Randomized.PIT.polyTimeComputable_sI, Randomized.PIT.polyTimeComputable_scanRun, Randomized.PIT.runOk, Randomized.PIT.sI, Randomized.PIT.scanRun, Randomized.PIT.validW}
\difficulty{4}
The Boolean function that sends each bit string $x$ to the one-bit string containing the
validity result for $x$ is computable by a deterministic Turing machine in time bounded
by a polynomial in the input length.
\end{theorem}

\begin{theorem}[Polynomial-Time Computability of a Bit-String Transformation]
\lean{Randomized.PIT.polyTimeComputable_runVal}
\label{Randomized.PIT.polyTimeComputable_runVal}
\leanok
\uses{Complexity.BinArith.getW, Complexity.BinArith.polyTimeComputable_getW, Complexity.PolyTimeComputable, Complexity.PolyTimeComputable.comp, Complexity.PolyTimeComputable.pairEncode, Complexity.polyTimeComputable_tail, Randomized.PIT.polyTimeComputable_sI, Randomized.PIT.polyTimeComputable_sL, Randomized.PIT.runVal, Randomized.PIT.sI, Randomized.PIT.sL, Turing.pairEncode}
\difficulty{4}
Let $R$ be the bit-string transformation that takes a string $D$, extracts the first and
second components from its nested pair encoding, drops the first bit of the resulting
first component, pairs what remains with the resulting second component, iterates the
second-component projection for a number of steps equal to one more than the paired
string’s length, and then extracts the first component of the resulting nested pair.
There is a deterministic Turing machine with finitely many states and work tapes that
computes $R(D)$ on every finite bit string $D$ in at most $C(|D|+1)^c$ steps, for some
natural numbers $C$ and $c$.
\end{theorem}

\begin{theorem}[Polynomial-Time Modular Evaluation]
\lean{Randomized.PIT.polyTimeComputable_evalModW}
\label{Randomized.PIT.polyTimeComputable_evalModW}
\leanok
\uses{Complexity.PolyTimeComputable, Complexity.PolyTimeComputable.comp, Randomized.PIT.evalModW, Randomized.PIT.gateRun, Randomized.PIT.polyTimeComputable_gateRun, Randomized.PIT.polyTimeComputable_runVal, Randomized.PIT.runVal}
\difficulty{3}
The function that takes a bit string $z$, initializes the encoded gate-evaluation state
from $z$, iterates the gate-state update for the prescribed number of rounds, and
returns the resulting value modulo the encoded modulus is polynomial-time computable:
there are a deterministic Turing machine with finitely many states and constants
$C,c\in\mathbb{N}$ that, on every input of length $n$, halts within $C(n+1)^c$ steps and
outputs this function's value.
\end{theorem}
